\section{Readouts without Enumerated Sub-Tasks}
\label{app:unspecified}
The controlled evaluation names the component predicates in the instruction. The examples below instead read intermediate operations from familiar programs whose component steps are not separately enumerated. They are qualitative illustrations from the exploratory slides, not part of the 253-specification evaluation.

\begin{figure}[ht]\centering
\includegraphics[width=\linewidth]{Figures/goalscope/qualitative/merge-sort.png}
\caption{\textbf{Reading intent while writing code.} At the marked positions, GoalScope reads descriptions of the operations the model is about to implement: (a) checking for an unmatched man, (b) finding the next free man, and (c) proposing to the next woman on his preference list. These descriptions express an operation rather than a particular token sequence; for example, finding an unmatched man can be implemented with different loop forms. The figure illustrates this distinction and does not itself compare readouts across alternative implementations. Later code is shown for context.}
\label{fig:freeform-matching}
\end{figure}

\begin{figure}[ht]\centering
\includegraphics[width=\linewidth]{Figures/goalscope/qualitative/stable-marriage.png}
\caption{\textbf{Reading intent while writing code.} At the marked positions, GoalScope reads descriptions of the operations the model is about to implement: (a) checking for an unmatched man, (b) finding the next free man, and (c) proposing to the next woman on his preference list. These descriptions express an operation rather than a particular token sequence; for example, finding an unmatched man can be implemented with different loop forms. The figure illustrates this distinction and does not itself compare readouts across alternative implementations. Later code is shown for context.}
\label{fig:freeform-matching}
\end{figure}

% \begin{figure}[ht]\centering
% \includegraphics[width=.64\linewidth]{Figures/slides/freeform_atoi.pdf}
% \caption{\textbf{A free-form string-to-integer example.} The highlighted source position is the indentation before \texttt{s = s.strip()}. A separate generic probe produces the displayed description about leading whitespace and signs. The later program is shown for context, but a cache at the highlighted token cannot depend on that future text. The description is a qualitative readout and need not match the exact granularity of the immediately following statement.}
% \label{fig:freeform-atoi}
% \end{figure}
